\CharacterSheet{example-character}{Título da personagem}{
% Acrescenta a personagem a content/characters/index.tex.
% Usa \Clue, \Space e \Character para fazer referências.
\shortdescription{Esta é uma breve descrição pública da personagem.}

\ObjectivesHeading{Objetivos}
\begin{itemize}
  \item \mainobjective{Este é um objetivo que a personagem tem de cumprir ao longo do jogo.}
  \item Este é um objetivo opcional que a personagem pode tentar cumprir ao longo do jogo.
\end{itemize}

\LosingConditionsHeading{Condições de derrota}
Aqui são indicadas as condições de derrota desta personagem. Se um jogador perder, os seus pontos não contam no fim do jogo.

\SecretsHeading{Segredos}
\begin{itemize}
  \item Este é um segredo que a personagem deve tentar esconder e que as outras querem descobrir.
\end{itemize}

\ActionsHeading{Ações}
\begin{itemize}
  \item Estas são ações que o jogador deve realizar para ajudar à interpretação da personagem.
\end{itemize}

\BackgroundHeading{Contexto}
Este é o passado da personagem, tal como será entregue ao jogador antes do início do jogo.

\GMNotes{Estas são notas sobre a personagem para o MJ consultar ao conceber o jogo. Que motivos tinha esta personagem? Que razões têm os outros jogadores para suspeitar dela? E que provas existem da sua inocência?}
}
